% Text measures shared by all five studies: the eight fixed word lists, how a script applied them
% without seeing the label or the rating, the two text checks run in every study, and the status of
% each use of the lists. Numbers come only from macros in generated/numbers.tex.
\label{sec:suppA-text}

In all five studies, a required text box for each rating asked participants to explain their reasoning for it, and every box came after the ratings it explained. This section describes how those explanations were measured, which was the same in every study: eight fixed word lists, applied by a script that never saw the label or the rating. It then describes the two checks on the explanations that every registration names, and the status of each use of the word lists. A registration is a dated record of a study's hypotheses and analysis plan (Section~\ref{sec:suppA-registrations}). Supplementary Note~7 gives what the explanations referred to, and each study's section gives its results on the two checks. Every explanation was written after its rating, so no result from the explanations shows why anyone rated as they did.

\subsection*{The eight word lists}

A word list here is a rule: a written list of words and phrases, in their ordinary forms, that a computer searches for in every explanation in the same way, after lowering its case. When an explanation contains a phrase on a list, the list flags that explanation. Supplementary Table~\ref{tab:suppA-wordlists} gives each list, what it flags and examples of what it matches.

% The eight fixed word lists applied to the written explanations in all five studies. Summarized
% from the rules file each study's text notebook writes (tables/text_rules.txt); no numbers.
% See supplement/tables/README_suppA.md.
\begin{table}[tbp]
\centering
\small
\caption{\textbf{The eight word lists.} Each list is a fixed search pattern that a script applies in the same way to every explanation, after lowering its case, without seeing the label or the rating. The examples are a selection; each study's rules file gives every pattern in full, with the phrases each list sets aside because they mean something else in these explanations. A list records a topic, whatever the writer said about it: `barely any detail' is detail talk.}
\label{tab:suppA-wordlists}
\begin{tabular}{@{}>{\raggedright\arraybackslash}p{2.6cm}>{\raggedright\arraybackslash}p{3.4cm}>{\raggedright\arraybackslash}p{\dimexpr\textwidth-9.4cm-6\tabcolsep\relax}>{\raggedright\arraybackslash}p{3.4cm}@{}}
\toprule
List (name in the code) & What it flags & Examples of what it matches & Further reading of a match\\
\midrule
Visible detail (\nolinkurl{says_detail}) & Detail and how much is in the picture & detail, intricate, complex, texture, shading, brushwork, layers, `so many different elements' & Asserted or denied\\
\addlinespace
Time (\nolinkurl{says_time}) & Time spent, duration or speed & time, hours, minutes, days, tedious, painstaking, spent, `took hours', quick, instant & Asserted or denied\\
\addlinespace
Skill (\nolinkurl{says_skill}) & Skill, craft and training & skill, craft, technique, talent, trained, mastery, expertise, professional, execution, `by hand' & Asserted or denied\\
\addlinespace
Concept (\nolinkurl{says_concept}) & The idea, meaning or thought behind the work & idea, concept, meaning, message, symbol, theme, story, represents, conveys, intention, `thought behind' & Asserted or denied\\
\addlinespace
Prompting or machines (\nolinkurl{says_prompt}) & Prompting, typing or generating, or a machine named in any role & prompt, typed, generated, input, AI, artificial intelligence, computer, algorithm, software, an image generator named by kind or brand & None\\
\addlinespace
Prior exposure (\nolinkurl{prior_exposure}) & Says the work looks familiar, already seen, generic or derivative & `seen before', familiar, `reminds me of', generic, cliché, typical, unoriginal, derivative, `nothing new' & Three values: familiarity asserted, familiarity denied, or neither\\
\addlinespace
Scale echo (\nolinkurl{anchor_echo}) & Repeats the rating scale's own wording & `moderate creativity', `a great deal of effort', `effort involved' & None\\
\addlinespace
Machine as maker (\nolinkurl{says_machine_made}) & Names a machine, program or image generator as making the work & `AI-generated', `machine-made', `created by a computer'; a person using a tool (`Alex used an AI') is set aside & Asserted, or not asserted (denied or only supposed)\\
\bottomrule
\end{tabular}
\end{table}


A list records a topic, not an endorsement: an explanation saying that a picture has barely any detail is talk about detail. For four lists (visible detail, time, skill and concept), a second reading records whether the matched phrase was asserted or denied. It looks for a negating word, such as `not', `never', `no', `little' or `lacks', in the thirty characters before the phrase, cut back to the start of the phrase's clause. The prior-exposure list uses the same window to separate a claim that a work looks familiar from a denial that it does. The machine-as-maker list counts a phrase as not asserted when it is denied or sits inside a supposition, such as one opened by `if', `whether' or `unless'; a hedged impression such as `looks AI generated' counts as asserted.

The lists, the negating words and every reading are written out in full in the text module that all five studies share (\nolinkurl{shared/python/study_text.py} in the analysis code; main text, Code availability). Each study's run also writes the rules in full to that study's \nolinkurl{tables/text_rules.txt}, with a SHA-256 checksum of that file, a fingerprint of its exact bytes that changes if one character changes, in \nolinkurl{tables/text_rules.sha256}. A third file, \nolinkurl{tables/text_rules_grounding.txt}, states what each list is meant to measure and the research it draws on. The five studies' rules files differ only in their title line and in the line that gives the number of explanations. A test in the shared folder (\nolinkurl{test_reproduces_study1.py}) runs the shared module on Study~1's explanations and confirms that it reproduces the outputs of Study~1's own text analysis exactly.

\subsection*{How the lists were written and applied}

We wrote the lists by reading Study~1's explanations with the label and the rating hidden. The lists were then frozen before any quantity that links them to the label or to the rating was computed, with four exceptions settled after results by label had been read. In September 2026 we settled how the prior-exposure list reads a negation, and wrote the machine-as-maker list itself. In October 2026 we settled that list's reading of whether a claim is asserted, and the rule that reads `little' as a denial only when it names an amount (`little detail') and not a size (`the little details'). Study~1's long registration records all four, and the paper's results that use those parts should be weighed with that in mind. The registrations of Studies 2 to 5 state that the lists run unchanged, and they did.

The script that applies the lists receives only identifiers, for the participant, the image and the rating asked about, and the text of each explanation. A check in the code stops the script if it receives any column that names the label, the rating or a background measure such as age. The label and the rating are joined to the flags only after the lists have run, and a second check stops the join if either had arrived earlier.

Our judgment entered these measures at one point: we read explanations to write the lists, as described above. No person and no language model labeled, coded or judged any single explanation for these measures, and no one read a flagged explanation to decide what it meant. Every flag and every count comes from the script, so each can be recomputed from the rules and the text alone. Samples of the explanations that each list flagged and did not flag, with the matched phrase marked, are saved with the analysis outputs, to be shared with the data (main text, Data availability), so that a reader can judge each list; nothing reported depends on what that reader concludes. Before Study~1 was registered, a language model had coded its explanations in exploratory work that saw the label (Section~\ref{sec:suppA-registrations}); that work was retired, and nothing in the paper uses it. One pretrained model does appear: a sentence-embedding model (all-mpnet-base-v2), which turns each explanation into a list of numbers whose distances reflect differences in meaning.\cite{reimers2019sentencebert,sentencetransformers2021allmpnet} It is used in two exploratory analyses only, a check on the topic model within each study and the measure of how far explanations moved in Supplementary Note~7, and it gives no explanation a label.

\subsection*{The two checks run in every study}

Study~1's long registration fixed two checks to run unchanged in every study where the rating exists, and each later registration names them. Both are registered secondary: named in a registration, but outside the set of tests that it names as confirmatory. In Studies 1 and 3 they ran on the explanations of both ratings; in Studies 2, 4 and 5 they ran on the creativity explanations only, as those registrations specify.

The first check asks how closely the words alone track the rating, and how well they tell the two label groups apart. Each explanation is reduced to counts of its content words, with function words such as `the' and `of' removed. Ridge regression, a regression that shrinks each word's weight toward zero so that more words than explanations can be fitted without memorizing them, predicts the rating, and its counterpart for a yes-or-no outcome predicts the label. Every prediction is made for participants left out of the fit: the participants are split into ten groups, and each group's explanations are predicted by a model fitted to the other nine. For the rating, the check reports the correlation between predicted and observed ratings. It sets that correlation beside the same correlation computed 200 times after each person's ratings were shuffled among that person's own explanations, which keeps each person's average and scrambles only which image got which rating.

For the label, the check reports the AUC, the chance that a randomly chosen AI-label explanation scores above a randomly chosen Human-label one, so that 0.5 is chance; its baseline shuffles the label across participants. Recovery is read as present when the estimate exceeds 95\% of the shuffled values and its 95\% confidence interval (CI), the range of values the data do not rule out at the 95\% level, clears chance. Otherwise recovery is not shown, which is not evidence that the words carry nothing. The check also reports how much the words add to the rating in telling the groups apart, as the AUC of the rating and the words together minus that of the rating alone. The words add to the rating when that interval lies wholly above zero, and the result is inconclusive when it includes zero. This first check is a description: it carries no prediction and is corrected with no other test.

The second check is the text counterpart of Study~1's H4, which predicted that under the AI label ratings would spread out more between people and less across one person's ratings of different images. It predicts that AI-label participants reuse their own words across their explanations more than Human-label participants do. For each participant it takes the average overlap among that person's own explanations, the share of content words that two explanations have in common. It subtracts the average overlap between that person's explanations and other participants' explanations of the same image, or of the same second work, which removes what any two people writing about the same picture share. The two label groups are then compared with explanation length held constant, with an interval from resampling participants within each group 2,000 times. Within each study the ratings checked are corrected together by Holm's method.\cite{holm1979simple} The reading is `supported' when the corrected $p$-value is below .05 with the AI group higher, `opposite direction' when it is below .05 with the AI group lower, and `inconclusive' otherwise. A $p$-value is the chance of a difference at least this large if there were none. An inconclusive reading is not evidence that the groups reuse their words equally.

Across the five studies, the words recovered the rating in every study and on every rating checked. They told the label groups apart in Study~1 on both ratings, in Study~3 on effort and, barely, in Study~5. They were not shown to do so in Studies 2 and 4 or on Study~3's creativity explanations, where the intervals reached chance. The words added to the rating in telling the groups apart in Studies 1 and 2, and the result was inconclusive in Studies 3, 4 and 5. Two of these readings turn on the third decimal place. Study~5's interval for the label begins just above chance, and those of Studies 2 and 4 just below it, although all three print as \STwoTextRecArmCreatLo{} at two decimals. Study~3's interval for what the words add on effort begins just below zero, although it prints as \SThreeTextRecIncrEffLo. Own-word reuse was greater under the AI label, as predicted, only in Study~1's effort explanations; every other reading was inconclusive. Supplementary Table~\ref{tab:supp-verdicts} summarizes these readings, and each study's section gives its estimates.

\subsection*{Other uses of the word lists, and their status}

Every other use of the lists is listed here with its status, so that no result from the explanations is read above it.

\emph{Study~1's five word-meaning tests} asked what `effort' and `creativity' referred to in the explanations. The registration names them as one exploratory set, corrected by the Benjamini--Hochberg method, which limits the expected share of false positives among results called significant.\cite{benjamini1995controlling} Their corrected results are descriptive readings, not verdicts (Supplementary Note~7; Supplementary Table~\ref{tab:suppA-denotation}).

\emph{Study~1's belief check} counts Human-label participants who mentioned machines or prompting, using the prompting list, and those who wrote that a work is or looks machine-made, using the machine-as-maker list read for assertion. The two counts, and a refit of the penalty without the second group, are registered secondary; the comparison with randomly chosen groups is exploratory description (Supplementary Note~2). Neither count measures disbelief in the label.

\emph{Study~1's familiarity association}, between calling an image familiar and its creativity rating within a person, is registered secondary description (Section~\ref{sec:suppA-study1}).

\emph{Study~4's refit} without the \SFourMachineMadeN{} of \SFourMachineMadeNHuman{} Human-label participants whom the machine-as-maker list flagged for writing that a machine made the work is registered secondary (Supplementary Table~\ref{tab:supp-verdicts}). \emph{Study~5's count} of such participants, \SFiveMachineMadeN{} of \SFiveMachineMadeNHuman, is exploratory.

\emph{Comparisons of how often the lists fired} under the two labels, in Study~3's explanations of originals and in the second-work explanations of Studies 2, 4 and 5, are exploratory and uncorrected (Supplementary Note~7).

\emph{The general description} of the explanations, run in every study, uses no word list. It scores the tone and concreteness of the words against published norms, the average ratings that large samples of people gave to single English words.\cite{warriner2013norms,brysbaert2014concreteness} It finds topics, groups of words that tend to occur together, with a model given no label and no rating, and it reports which words a lasso, a regression that sets most word weights to exactly zero, keeps when it predicts the label or the rating. Every registration names this description as exploratory, and it carries no verdict.
